\answer{ex:06:1}{$e=(1-e^{-T_a/\tau_e})e^{-d/\tau_e}$. (а)~$0.110$ против $4.5\cdot10^{-5}$, в $\approx2.5\cdot10^3$ раз. (б)~$\tau_e^*=T_a/\ln(1+T_a/d)=0.59\,\text{с}$.}
\answer{ex:06:2}{Скорость $\nu_x\nu_y(A_+\tau_+-A_-\tau_-)=0.8A_+$ в секунду, $\approx2900A_+$ за час; $\tau_-=\tau_+/0.6=33\ms$.}
\answer{ex:06:3}{$\mathbf e_1$ при $\mu^2+s_1^2>s_2^2$; здесь $1.01>0.25$, и нейрон выучит $\mathbf e_1$, хотя разброс вдоль $\mathbf e_2$ в $25$ раз больше: правило находит главный вектор матрицы корреляций, а не ковариаций.}
\answer{ex:06:4}{$V=Re^{-(t_c+L-t)/\tau_\gamma}$ между лампочкой и соком, так что $\dot V=V/\tau_\gamma$; всплеск площади $Re^{-L/\tau_\gamma}$: $0.61R$ и $0.78R$.}
\answer{ex:06:5}{Отношение сигнала к шуму $1/(N+1)$, попыток $\propto N+1$: $5\cdot10^5$ попыток $\approx1$~мес. и $5\cdot10^{11}$ попыток $\approx8\cdot10^4$~лет.}
\answer{ex:06:6}{(а)~$k_2=1/\tau_d$, $k_1=1/\tau_d-1/\tau_\gamma$. (б)~Ложная ошибка $-(V/\tau_\gamma)\,\varepsilon/(1+\varepsilon)$, $\varepsilon=\tau_d/\tau_\gamma$, откуда $\tau_d\le0.105\,\text{с}$.}
\answer{ex:06:7}{$0.46$, $0.90$, $0.99$, $0.95$ и $0.70$ при $d=3\,\text{с}$. Точки ложатся на одну кривую от доли следа $F=(1-e^{-T_{\text{cue}}/\tau_e})e^{-d/\tau_e}$: обучение при $F\gtrsim0.1$, случайный выбор при $F<10^{-2}$.}
\answer{ex:06:8}{$\eta^*=1/\bigl(2\sigma^2(N+2)\bigr)$, за $N+2$ попытки; при $m$ дрожащих из $N$ --- за $N(m+2)/m\ge N+2$: разреженность не помогает.}
